\section{Finite computations and the main lower bound}\label{auto:section}

We complete the concrete obligations deferred in
Section~\ref{sec:diagrams}: constructing the small source automaton and
proving the finite-computation representation. The order-reversing map
has already been derived in Proposition~\ref{auto:order-reversal};
we then combine it with Theorem~\ref{amp:image} to finish the main proof.
Throughout, we use the automaton model fixed in the introduction,
including stay moves, finite acceptance, and acceptance in the initial
configuration.

\subsection{The source automaton}

Fix a set $H$ of size $h\geq2$. Recall that the alphabet is
$\Sigma_H=\mathcal R_H$, the word product is $r(w)$ with
$r(\varepsilon)=I_H$, and
\begin{equation}\label{auto:language}
 L_H=\{w\in\Sigma_H^*:r(w)\ne\varnothing\}.
\end{equation}

\begin{lemma}\label{auto:source}
The language $L_H$ is recognized by a 2NFA with exactly $h+2$ states.
\end{lemma}
\begin{proof}
Use one initial state $q_{\mathrm{in}}$, the $h$ states in $H$, and one
accepting state $q_{\mathrm{acc}}$, all distinct.  From
$q_{\mathrm{in}}$ on the left endmarker, the machine moves right into any
chosen state of $H$.  On a letter $R\in\mathcal R_H$, it may move right
from $p\in H$ to $q\in H$ precisely when $(p,q)\in R$.  At the right
endmarker, every state in $H$ has a stay transition to
$q_{\mathrm{acc}}$.  There are no other transitions, and
$q_{\mathrm{acc}}$ is the only accepting state.

On $R_1\cdots R_k$, an accepting computation is exactly a sequence
$p_0,\ldots,p_k\in H$ with $(p_{i-1},p_i)\in R_i$ for
$1\leq i\leq k$.  Such a sequence exists exactly when the relation
product is nonempty.  For $k=0$, the initial move reaches the right
endmarker directly, and the machine accepts.  This agrees with
$r(\varepsilon)=I_H\ne\varnothing$.
\end{proof}

\subsection{Finite computations as diagram paths}
\label{auto:finite-computations}

We now prove the representation promised in
Lemma~\ref{auto:representation}. The construction records arbitrary
finite computations, including stay moves and repeated crossings of the
same cut. It does not require termination of every computation.

\begin{proof}[Proof of Lemma~\ref{auto:representation}]
Let $B$ be the $s$-state automaton in the lemma, with state set $Q$ and
initial state $q_0$. Introduce
one new state $q_\dagger$.  From every original accepting state, at
every cell, add a stay transition to $q_\dagger$.  In state
$q_\dagger$, move right regardless of the scanned symbol, finally
exiting past the right endmarker in state $q_\dagger$.  This last exit
is only a device for representing acceptance, not a transition of the
original machine.  All original transitions remain subject to the
endmarker restrictions.

A finite augmented computation from the original initial configuration
to the designated exit exists if and only if $B$ accepts.  Indeed, an
accepting original computation can be followed by the added stay and
rightward sweep.  Conversely, the first entry into $q_\dagger$ must
come from an original accepting configuration.  This argument includes
the case that $q_0$ itself is accepting.  Infinite computations create
no additional finite successful computation.

Put $\widehat Q=Q\cup\{q_\dagger\}$ and $m=s+1$.  Take the
directional label sets to be two disjoint copies
\[
  D^+=\{q^+:q\in\widehat Q\},
  \qquad
  D^-=\{q^-:q\in\widehat Q\}.
\]
For a cell symbol $c$, including either endmarker with its boundary
rules, let $E_c$ be the relation on $\widehat Q$ given by the augmented
stay transitions.  Let $T_c^+$ and $T_c^-$ be the relations given by
the augmented rightward and leftward transitions, respectively; their second
coordinates record the state \emph{after} the move.  At the right
endmarker, $T_c^+$ includes only the designated augmented exit, and at
the left endmarker $T_c^-$ is empty.  Define a diagram $d_c$ by
\begin{align*}
 (q^+,p^+)\in F_{d_c}
 &\quad\Longleftrightarrow\quad
 (q^-,p^+)\in R_{d_c}
 \quad\Longleftrightarrow\quad
 (q,p)\in E_c^*T_c^+,\\
 (q^-,p^-)\in B_{d_c}
 &\quad\Longleftrightarrow\quad
 (q^+,p^-)\in L_{d_c}
 \quad\Longleftrightarrow\quad
 (q,p)\in E_c^*T_c^-.
\end{align*}
Thus each local diagram edge represents a finite sequence of stays,
possibly empty, followed by one move out of the cell.  The incoming
sign specifies the side of entry, not additional machine memory; the
transition rules themselves depend only on the state and symbol.

Set $\lambda=d_{\mathrm{left}}$ and
$\rho=d_{\mathrm{right}}$ for the two endmarkers, and set
\[
  \tau(c_1\cdots c_k)=d_{c_1}\cdots d_{c_k},
  \qquad \tau(\varepsilon)=1_{\mathcal T_m}.
\]
This is a monoid homomorphism.  Every finite augmented computation
ending in the designated exit decomposes into successive visits to
individual cells, each consisting of finitely many stays and one
exit move.  These visits give a path through the product of the cell
diagrams.  In the opposite direction, each local edge in a finite
diagram path has a finite transition witness.  Expanding these
witnesses and concatenating them gives an augmented computation,
because the exit state from one cell is exactly the entry state to
the next.  Neither direction assumes that a cut is crossed at most
once.

The fictitious entry into the left endmarker with label $q_0^+$
represents the initial configuration; it is not an additional machine
transition.  No path can leave the tape to the left, and a path can
leave to the right only through the designated exit.  Hence
\eqref{auto:recognition} holds with $a=q_0^+$ and $b=q_\dagger^+$.
For the empty word, the two endmarker cells are adjacent, and the
identity $\tau(\varepsilon)$ gives precisely their product.  A stay
cycle causes no difficulty: only its finite traversals are represented
in $E_c^*$.
\end{proof}

\begin{remark}[Other standard conventions]\label{auto:conventions}
The exponential conclusion also survives the usual changes of starting
position or acceptance convention, with possible changes to the additive
state offsets. To use our representation for a machine whose convention
requires a positive transition into an accepting state, add a fresh
nonaccepting initial copy with the same outgoing transitions, keeping all original states and
transition destinations. This prevents the initial configuration alone
from accepting, while preserving every acceptance after a transition.
Starting on the first input cell can be simulated from the left
endmarker by one additional initial state. If success must occur at a designated marker or on a specified
exit, enable the added success routine only at that local accepting
event. These versions of the path representation use $m=s+O(1)$ labels
per sign, and the source automaton still uses $h+O(1)$ states. The exact
offsets in Theorem~\ref{main:theorem} refer to the convention fixed in
the introduction.
\end{remark}

\begin{proof}[Proof of Theorem~\ref{main:theorem}]
For any $n\geq4$, let $h=n-2$, choose a set $H$ of size $h$, and use
the $n$-state source automaton from Lemma~\ref{auto:source}.
If an $s$-state 2NFA recognizes its complement, then
Proposition~\ref{auto:order-reversal} supplies the order-reversing
relation image required by Theorem~\ref{amp:image}.  With $m=s+1$,
that bound gives
\[
  2(s+1)\geq 2^{\lfloor(h-2)/127\rfloor}
            =2^{\lfloor(n-4)/127\rfloor}.
\]
Rearranging proves the stated lower bound.  Since its right-hand side
grows exponentially with $n$, no fixed polynomial can bound the
complementation cost for all finite alphabets.
\end{proof}

\subsection{The determinization consequence}

The lower bound also applies indirectly to deterministic simulation.
Here we use one external automata transformation: deterministic two-way
automata can be complemented with linear state overhead independent of
the alphabet, even when the original computation may fail to halt
\cite{GeffertMereghettiPighizzini2007}. The possible nonaccepting loops
are the reason that exchanging accepting and rejecting states alone
does not give this transformation.

\begin{corollary}\label{auto:determinization}
There are absolute constants $c>0$ and $n_0$ such that, for every
$n\ge n_0$, every 2DFA recognizing the language of the $n$-state 2NFA
$A_n$ in Theorem~\ref{main:theorem} has at least
\[
 c\,2^{\lfloor(n-4)/127\rfloor}
\]
states. Hence no polynomial state bound independent of the finite
alphabet can determinize all 2NFAs.
\end{corollary}
\begin{proof}
We use only the linear-overhead conclusion of
\cite{GeffertMereghettiPighizzini2007}, allowing constant-factor changes
for model conventions. The standard marker model and its normalization
are described in the precursor \cite[Section~2 and Lemma~3.1]{GeffertMereghettiPighizzini2005}.
Stay moves can be replaced by a two-step
excursion that stores the destination state and the return direction
in a fixed number of copies of the state set. The excursion goes left
at the right endmarker and right elsewhere. Even on the empty word,
the two distinct endmarker cells provide the needed adjacent cell.
Auxiliary states are nonaccepting, and the destination state is entered
only on returning to the original cell. A reached accepting
configuration, including the initial configuration, can instead start
a fixed marker sweep to a designated accepting halt. Missing
transitions can halt and reject. These changes require $O(s+1)$ states
for an $s$-state 2DFA and preserve finite acceptance. The cited
transformation supplies the substantive additional step of handling
infinite nonaccepting computations.

Thus an $s$-state deterministic recognizer for $L(A_n)$ has a
deterministic complement, also a 2NFA, with at most $C(s+1)$ states
for an absolute constant $C$. Theorem~\ref{main:theorem} gives
\[
 C(s+1)\ge \tfrac12\,2^{\lfloor(n-4)/127\rfloor}-1.
\]
For all sufficiently large $n$ this implies the asserted bound with
an absolute $c>0$. Since the same alphabets are finite for every $n$,
it also excludes every polynomial bound uniform over alphabets.
\end{proof}

For comparison, the companion theorem
\cite[Theorem~1.1]{CompanionDeterminization2026} gives
\[
 4(s+2)^2\geq 2^{\lfloor(h-2)/31\rfloor}
\]
for an $s$-state deterministic recognizer of the same language $L_H$.
That theorem also states the improvement from $s+2$ to $s+1$ when the
initial configuration counts as acceptance, as it does here. Its
exponent therefore grows as $h/62$, compared with $h/127$ in
Corollary~\ref{auto:determinization}. The source constructions use
$h+3$ and $h+2$ states, respectively: the companion includes a rejecting
sink, which the partial source automaton above does not need. Its
matching-diagram proof gives a stronger deterministic bound independently
of the complementation theorem proved here.
